\chapter{Pre-Engagement}

\section{Overview}
This chapter highlights the Pre Engagement phase of the penetration testing.
\\

All the Rules of Engagement (RoE) for the penetration testing activities are defined in this section.

\section{Purpose and Objectives}

The purpose of this penetration testing activity is to evaluate the security
posture of the intentionally vulnerable virtual machine "Corrosion: 1".
\\

The penetration testing activities will be performed in a black-box fashion.
\\

The activity is conducted as part of the Penetration Testing and Ethical
Hacking exam and aims to demonstrate practical knowledge of offensive
security methodologies.
\\

The objectives of the assessment are:

\begin{itemize}
    \item Identify security vulnerabilities present within the target system.
    \item Perform reconnaissance and enumeration activities.
    \item Exploit discovered vulnerabilities where technically feasible.
    \item Demonstrate penetration testing methodologies and techniques.
    \item Document discovered vulnerabilities, exploitation procedures,
    evidence, and remediation recommendations.
\end{itemize}

The penetration test is conducted exclusively for educational purposes within
an authorized laboratory environment.


\section{Authorization}

This section provides formal authorization for Mariano Aponte to perform
penetration testing activities against the specified target system.
\\

The target system is an intentionally vulnerable virtual machine obtained from
VulnHub and deployed inside an isolated virtual laboratory environment. The
system is designed for security training purposes and does not represent a
production asset.
\\

The tester is authorized to perform security testing activities against the
target machine within the scope and limitations described in this chapter.


\section{Scope of Testing}

\subsection{In-Scope Assets}

The following asset is explicitly authorized for testing:

\begin{itemize}
    \item \textbf{Target:} Corrosion: 1 virtual machine.
    \item \textbf{Environment:} Isolated virtual laboratory network.
\end{itemize}

All services, applications, configurations, user accounts, and files contained
inside the target virtual machine are considered within scope.


\subsection{Out-of-Scope Assets}

The following assets are explicitly excluded from the scope:

\begin{itemize}
    \item The physical host machine running the virtualization software.
    \item Other virtual machines present in the laboratory environment.
    \item External networks and Internet-accessible systems.
    \item Any system not explicitly identified as the target machine.
\end{itemize}

No testing activity shall intentionally affect systems outside the authorized
scope.


\section{Authorized Testing Activities}

The following activities are authorized against the target virtual machine:

\begin{itemize}
    \item Network discovery and host identification.
    \item Port scanning and service enumeration.
    \item Vulnerability identification and assessment.
    \item Web application testing.
    \item Configuration analysis.
    \item Credential testing against exposed services.
    \item Controlled exploitation of identified vulnerabilities.
    \item Privilege escalation attempts.
    \item Post-exploitation analysis limited exclusively to the target machine.
    \item Collection of technical evidence for reporting purposes.
\end{itemize}

Since the target machine is intentionally vulnerable and isolated, exploitation
activities are permitted when required to demonstrate penetration testing
techniques.


\section{Prohibited Activities}

The following actions are prohibited:

\begin{itemize}
    \item Attacking systems outside the defined scope.
    \item Performing denial-of-service attacks against external systems.
    \item Modifying or damaging the host operating system.
    \item Accessing unrelated personal or confidential information.
    \item Performing security tests against third-party infrastructure.
    \item Publishing zero-day vulnerabilities without authorization.
\end{itemize}


\section{Testing Methodology}

The penetration testing activity follows the methodology defined by the Penetration Testing Execution Standard (PTES). 
\\

PTES provides a structured framework for conducting penetration tests by dividing the
assessment process into seven distinct phases. The following phases define the
approach adopted during the assessment of the "Corrosion: 1" vulnerable
virtual machine.

\subsection{Pre-Engagement Interactions}

The pre-engagement phase establishes the objectives, scope, rules, and
limitations of the penetration test.
\\

During this phase, the following elements are defined:

\begin{itemize}
    \item Identification of the target system and testing environment.
    \item Definition of authorized activities and testing boundaries.
    \item Confirmation that the target is an intentionally vulnerable machine
    deployed in an isolated laboratory environment.
    \item Definition of reporting requirements and documentation procedures.
\end{itemize}

The Rules of Engagement document represents the authorization and scope
definition for the entire penetration testing activity.

\subsection{Intelligence Gathering}

The intelligence gathering phase consists of collecting information about the
target environment in order to identify potential attack vectors.
\\

Activities performed during this phase may include:

\begin{itemize}
    \item Passive information gathering where applicable.
    \item Identification of network information related to the target.
    \item Discovery of exposed services and applications.
    \item Collection of information useful for subsequent enumeration activities.
\end{itemize}

The objective of this phase is to build an initial understanding of the target's
attack surface.

\subsection{Threat Modeling}

The threat modeling phase evaluates possible attack scenarios against the
identified target.
\\

During this phase, potential threats are analyzed based on:

\begin{itemize}
    \item Exposed services and applications.
    \item Identified vulnerabilities.
    \item Possible attacker objectives.
    \item Potential impact of successful exploitation.
\end{itemize}

The analysis is used to prioritize testing activities and determine the most
relevant attack paths.

\subsection{Vulnerability Analysis}

The vulnerability analysis phase focuses on identifying weaknesses within the
target system.
\\

Activities include:

\begin{itemize}
    \item Service and version enumeration.
    \item Identification of known vulnerabilities.
    \item Manual analysis of discovered services.
    \item Assessment of vulnerability severity and exploitability.
\end{itemize}

The results of this phase determine which vulnerabilities will be validated
through controlled exploitation.

\subsection{Exploitation}

The exploitation phase consists of performing controlled attacks against
identified vulnerabilities in order to verify their practical impact.
\\

Authorized activities include:

\begin{itemize}
    \item Exploitation of vulnerable services and applications.
    \item Validation of discovered security weaknesses.
    \item Collection of proof-of-concept evidence.
    \item Documentation of successful attack paths.
\end{itemize}

All exploitation activities are limited exclusively to the \textit{Corrosion: 1}
virtual machine.

\subsection{Post-Exploitation}

The post-exploitation phase evaluates the impact of successful exploitation and
determines the level of access obtained.
\\

Activities may include:

\begin{itemize}
    \item Identification of obtained privileges.
    \item Assessment of privilege escalation opportunities.
    \item Evaluation of accessible system resources.
    \item Collection of evidence required for reporting.
\end{itemize}

Post-exploitation activities are performed only to demonstrate security impact
and do not involve unnecessary access or data collection.

\subsection{Reporting}

The reporting phase documents all activities, findings, and conclusions of the
penetration test.
\\

The final report includes:

\begin{itemize}
    \item Description of the testing methodology.
    \item Scope and limitations of the assessment.
    \item Tools and techniques employed.
    \item Identified vulnerabilities.
    \item Exploitation details and evidence.
    \item Risk assessment and impact analysis.
    \item Recommended remediation actions.
\end{itemize}

The final documentation is submitted to Prof. Arcangelo Castiglione for
evaluation as part of the Penetration Testing and Ethical Hacking exam.


\section{Data Handling}

Information collected during the penetration test shall be handled responsibly.
\\

The tester agrees to:

\begin{itemize}
    \item Use collected information only for academic evaluation purposes.
    \item Avoid unnecessary extraction of data.
    \item Preserve only evidence required for the final report.
    \item Remove temporary testing artifacts when possible.
\end{itemize}


\section{Risk Management}

Although the target system is intentionally vulnerable, penetration testing
activities may cause service interruption, system instability, or modification
of data contained within the virtual machine.
\\

The tester is responsible for maintaining the isolation of the laboratory
environment and ensuring that testing activities remain limited to the
authorized target.
\\

Virtual machine snapshots may be used to restore the environment if necessary.


\section{Reporting and Communication}

The final penetration testing report will contain:

\begin{itemize}
    \item Scope and methodology.
    \item Tools and techniques used.
    \item Identified vulnerabilities.
    \item Exploitation procedures.
    \item Evidence of successful exploitation.
    \item Security impact analysis.
    \item Recommended mitigation strategies.
\end{itemize}

The report will be submitted to Prof. Arcangelo Castiglione as part of the
Penetration Testing and Ethical Hacking examination.


\section{Duration of Testing}

Testing activities are authorized during the exam period assigned by the
course supervisor.
\\

The tester may perform penetration testing activities against the target
virtual machine during this period, provided that all rules defined in this
document are followed.


\section{Declaration of Authorization}

By accepting these Rules of Engagement, the involved parties
acknowledge that:

\begin{itemize}
    \item The penetration testing activity is authorized.
    \item The target system is intentionally vulnerable and used exclusively
    for educational purposes.
    \item Testing activities are restricted to the defined scope.
    \item The tester agrees to follow ethical hacking principles.
\end{itemize}